% Experimental protocol.

\section{Experimental protocol}\label{sec:protocol}

\subsection{The cases}\label{sec:cases}

A case is an optical setup, a device, a target (a list of traps or a shaped potential), an aberration, a method set, a seed and, for dynamic cases, a particle and a rearrangement. There are \nCases{} cases in four categories (Table~\ref{tab:cases}): arrays (a $3\times3$ grid, a $10\times10$ grid with camera feedback, a $32\times32$ array of \maxTraps{} traps, the calibration task of the analytic and Gerchberg--Saxton baselines, and a disordered array of \nDisordered{} traps on integer bins at least five bins apart), shaped and three-dimensional targets (a disc potential with the crosstalk model, a three-plane array with vortices, a vortex ring with a line, and a three-plane volume at the geometry of the released learned model), device and optics (high numerical aperture, a device-limited array, an aberrated array with and without correction) and dynamics and controls (an atom-rearrangement sequence and a trap outside the field, which is refused by the first contract and never computed). The baked cases span \minTraps{} to \maxTraps{} traps. The optical setups are listed in Table~\ref{tab:setups}; the bead setup is in water, and the setup of C14 is the volume of the released learned model.

\begin{table}[t]\centering\footnotesize
\caption{The fourteen cases: category, optical setup, target and the methods that ran in the bake. C11 is refused by the input contract and never computed.}\label{tab:cases}
\begin{tabular}{@{}l l l l r >{\raggedright\arraybackslash}p{0.14\textwidth}@{}}\toprule Case & Category & Setup & Target & Methods & Reference\\\midrule
C01 & Arrays & 850 nm, NA 0.5 & 9 traps & 14 & weighted GS\\
C02 & Arrays & 850 nm, NA 0.5 & 100 traps & 14 & phase-fixed weighted GS\\
C03 & Arrays & 850 nm, NA 0.5 & 1024 traps & 7 & phase-fixed weighted GS\\
C04 & Shaped and 3D & 850 nm, NA 0.5 & disc potential, 8-bit & 6 & device-aware CG with pixel crosstalk\\
C05 & Shaped and 3D & 1064 nm, NA 1.2 & 51 traps & 10 & direct search\\
C06 & Shaped and 3D & 1064 nm, NA 1.2 & 25 traps & 8 & phase-fixed weighted GS\\
C07 & Device and optics & 1064 nm, NA 1.2 & 25 traps & 13 & phase-fixed weighted GS\\
C08 & Device and optics & 850 nm, NA 0.5 & 64 traps, 4-bit, fill 0.85 & 7 & device-aware weighted GS\\
C09 & Device and optics & 850 nm, NA 0.5 & 16 traps, aberrated & 6 & Zernike correction\\
C10 & Dynamics and controls & 850 nm, NA 0.5 & 100 traps & 4 & Hungarian assignment with phase interpolation\\
C11 & Dynamics and controls & & one trap at 200 $\mu$m, outside the field & 0 & refused\\
C12 & Arrays & 850 nm, NA 0.5 & 100 traps & 7 & weighted GS\\
C13 & Arrays & 850 nm, NA 0.5 & 819 traps & 7 & phase-fixed weighted GS\\
C14 & Shaped and 3D & 850 nm, NA 0.5 & 48 traps & 8 & multi-plane CNN hologram\\
\bottomrule\end{tabular}\end{table}


Every case states numeric expected ranges for its metrics, measured on three seeds, and the bake fails outside them. All cases are synthetic: the parameters are stated and no measured sample is used. The cases C01, C02, C03, C12 and C13 sit on integer-bin lattices so that the Fourier form applies and the live lane keeps its time budget; the cases C05 to C10 keep off-bin traps because they carry depth, vortices, a high numerical aperture, a device, an aberration or a rearrangement, which the Fourier form cannot represent. A device case excludes traps off the focal plane, vortex charge and an aberration, because the device model is read at one focal plane. The reference method of a case is the one whose maps the replay shows, and the method the reader picks among those valid for the case is a variant, as is the device (ideal or the case's own).

\begin{table}[t]
\centering\footnotesize
\caption{The optical setups of the cases, read from the committed manifests and traces. The device row of C08 is the device the benchmark also uses for its device cells.}
\label{tab:setups}
\begin{tabular}{@{}l r r r r r r r r@{}}
\toprule
Setup & $\lambda$ (nm) & $f$ (mm) & NA & SLM pixels & pitch ($\mu$m) & pad & bin ($\mu$m) & half field ($\mu$m)\\
\midrule
Atom (C01--C04, C09--C13) & \atomLambda & \atomFocal & \atomNa & $\atomNpix^2$ & \atomPitch & \atomPad & \atomBin & \atomHalf\\
Bead (C05--C07, water) & \beadLambda & \beadFocal & \beadNa & $\beadNpix^2$ & \beadPitch & \beadPad & \beadBin & \beadHalf\\
C08 & \devcLambda & \devcFocal & \devcNa & $\devcNpix^2$ & \devcPitch & \devcPad & \devcBin & \devcHalf\\
C14 & \voluLambda & \voluFocal & \voluNa & $\voluNpix^2$ & \voluPitch & \voluPad & \voluBin & \voluHalf\\
\bottomrule
\end{tabular}

\smallskip
\footnotesize The bead setup is in water of index \beadNimm; C14 is the volume of the released learned model. C08 device: \devcBits-bit phase, crosstalk $\sigma_K=\devcXtSigma$, $q=\devcXtQ$ on $3\times3$ sub-pixels, fill factor \devcFill, flicker \devcFlicker{} rad.
\end{table}

\subsection{Metrics, baselines and acceptance criteria}\label{sec:protocol-metrics}

The metrics are those of Eq.~\eqref{eq:metrics}: the efficiency, the uniformity and the relative standard deviation of the trap intensities, plus the zero-order and ghost fractions of the power budget, all computed by the scorer from the final phase through the case's device. For a shaped target the two required slots carry the signal-region metrics: the relative-standard-deviation slot holds the root-mean-square fractional error of Eq.~\eqref{eq:m08-err} in percent after the scale fit, and the uniformity slot holds one minus its fraction, clipped at zero. For the learned volume network the metric is the accuracy of Eq.~\eqref{eq:m17}, and the comparison is the median per-volume ratio of the accuracies at equal wall-clock.

The baselines are run by the product on the same task, never only quoted: the analytic and Gerchberg--Saxton methods M01 to M06 are the baselines of the weighted methods; GS is the baseline of the optimization methods at equal iterations; the volumetric Gaussian GS is the baseline of the learned network at equal wall-clock; the phase-fixed weighted GS is the baseline of every proposed method; and the random start is the baseline of the warm start. The published values that a method's acceptance criterion quotes are used as anchors, and the report states where the engine reproduces them and where it does not. The acceptance criterion of each method is one measurable statement with one test (Table~\ref{tab:ladder}); the measured value of each, with its seeds, is Table~\ref{tab:acceptance}. A criterion is confirmed when the test passes, and a method that fails its criterion on DinHot's grids is reported as measured, not hidden: the criterion of the warm start is refuted, the criterion of the aberration-robust hologram was narrowed after a measurement, and the criterion of the freeze index of the phase-fixed method is met only with a larger value than the source's.

\begin{longtable}{@{}>{\raggedright\arraybackslash}p{0.34\textwidth} >{\raggedright\arraybackslash}p{0.58\textwidth}@{}}
\caption{The methods of the studio and the acceptance criterion the design document fixes for each: the published methods and the methods proposed here.}\label{tab:ladder}\\
\toprule Method & Acceptance criterion\\\midrule\endfirsthead
\toprule Method & Acceptance criterion\\\midrule\endhead
\multicolumn{2}{@{}l}{\textit{Published methods}}\\
Superposition of prisms and lenses & \footnotesize on the 10x10 task with a flat beam, efficiency within 0.10 of 0.29 and uniformity below 0.05\\
Random superposition & \footnotesize on the 10x10 task with a flat beam, efficiency above superposition of prisms and lenses and at least 0.65; uniformity below 0.35 (DinHot measured 0.25 at 256 squared against 0.01 at 768 squared in the source; the ordering, not the value, transfers)\\
Random mask & \footnotesize efficiency about 1/M; at the source's 768 squared grid, uniformity within the spread of 0.58\\
Gerchberg-Saxton, point-trap and FFT forms & \footnotesize the two forms agree to $10^{-9}$ rad; 30 iterations on the 10x10 task with a flat beam reach efficiency at least 0.9 and uniformity near 0.6\\
Generalized adaptive additive & \footnotesize uniformity above GS at equal iterations on the 10x10 task\\
Weighted GS & \footnotesize efficiency and uniformity at least 0.9 within a median of 15 iterations on the 10x10 task with a flat beam (the source reports about 10); the 3x3x3 lattice above 0.9 and 0.98\\
Phase-fixed weighted GS & \footnotesize fewer iterations than weighted GS to a relative standard deviation of 0.5 percent on a 50x30 array, with the phase fixed after 60 iterations in the acceptance run (the code default `n\_fix` is 12, as in Kim et al. 2019; at 12 it stalls at 3.58 percent on DinHot's grid, re-measured 2026-10-30); from the same start phase-fixed weighted GS matches slmsuite 0.5.0's WGS-Kim to 1e-4 percent (the first field is not reweighted, g\_0 = 1)\\
MRAF & \footnotesize on disc, ring and bar targets, signal-region error at least five times below GS at the cost of efficiency\\
Direct search & \footnotesize uniformity at least 0.98 on the 10x10 task; under 5 percent rms variation on a quasi-periodic set\\
Conjugate-gradient minimization & \footnotesize the analytic gradient matches finite differences; a square lattice below 1 percent error\\
Wirtinger gradient descent with Adam & \footnotesize lower trap-intensity error than GS at equal iterations on the 10x10 task\\
Closed-loop WGS on an emulated camera & \footnotesize from a 19 percent rms open-loop spread, at most 2.1 percent in 20 iterations at gain 0.7; the acceptance runs on the emulated camera (`core/camera.py`), which reproduces the published feedback curves (Nogrette 2014: 19 to 1.4 percent in 20 iterations; Kim 2019: 22 percent to 1.59 plus or minus 0.18 percent in three phase-fixed steps over 8 ensembles, 1.4 in the histogram of one run; the bounds 2.1 and 1.8 are set above them)\\
Device-aware CG with pixel crosstalk & \footnotesize on the crosstalk device, lower error than conjugate-gradient minimization without the model on a large disc; a staircase line search lets it run its budget on a quantizing device, where its scored spread does not improve systematically (a measured negative)\\
Zernike correction and sensorless modal recovery & \footnotesize recovers an injected set of Noll modes 4 to 11 within 5 percent of their amplitude\\
Laguerre-Gaussian vortex traps & \footnotesize the ring phase winds by 2 pi l with an error under 0.1 rad\\
Hungarian assignment and linear phase interpolation & \footnotesize no crossing moves for alpha 1.5 on a 7x7 to 3x3 sort; a moving trap above 70 percent of its mean intensity on the frames themselves (the minimum over the frames of the sequence over their mean; between frames, under the blend model of flicker-bounded phase-continuous sequence's transient, it is not: notebook 07-04)\\
Multi-plane CNN hologram with an optional volumetric GS refinement (DeepCGH-style) & \footnotesize on 162 held-out three-plane volumes, the hybrid (the network's phase, refined by volumetric 3D GS inside the same wall-clock budget only when that raises the volume accuracy) reaches a higher accuracy than volumetric 3D GS alone at that budget (measured median ratio 1.140); the refinement was kept on none of the volumes, so the measured method is the network alone\\
Optimal-transport initialized MRAF & \footnotesize lower error than MRAF from a random start on the shaped targets\\
Multi-plane GS & \footnotesize on a three-plane target, every plane mean within 25 percent of the overall mean (volumetric 3D GS of Whyte and Courtial)\\
\multicolumn{2}{@{}l}{\textit{Methods proposed here}}\\
Device-aware weighted GS & \footnotesize on the crosstalk device, relative standard deviation at least two times below phase-fixed weighted GS, efficiency at most 10 percent lower; measured against the corrected phase-fixed weighted GS and against slmsuite's WGS-Kim alike: a median of 25.8 times (15.9 to 33.3 over five seeds) at 0.995 of the efficiency, and 2.0 times on the full C08 device (the bake, three seeds; 2.2 in notebook 03-04) with 4-bit phase and flicker; in the benchmark, on the C08 device at 100 iterations with the default parameters, device-aware weighted GS meets both thresholds on 2 of 5 geometries (G2 ring of 36 traps, 2.28 times at 0.965 of the efficiency, G3 random 100 traps, 4.69 times at 0.938 of the efficiency); lowers the spread without the factor of two on G6 (1.73 times at 1.06); keeps its start on G1 (0.97), G4 (1.05); device-aware weighted GS's cells hold 2 seeds\\
Flicker-bounded phase-continuous sequence & \footnotesize minimum transient intensity of a moving trap at least 0.9 of its mean, against LPI on the same assignment: with the frames that keep the peak step per frame at most 0.2 Fourier units (a bound set on C10's sort) it holds on most seeds (acceptance task, 13 frames at 0.31 Fourier units per frame, seeds 1 to 3: flicker-bounded phase-continuous sequence 0.916 to 0.926, LPI 0.551 to 0.709; long lattice moves: median 0.933 and 0.936, 6 of 8 seeds, and flicker-bounded phase-continuous sequence above LPI on 7 of 8 seeds of each); on C10's long sort at 25 frames LPI leads in every seed, and flicker-bounded phase-continuous sequence needs the end phases and 137 frames to reach 0.902 to 0.917 on 5 of 6 seeds of the plan of the Hungarian assignment with phase interpolation (median 0.911; LPI 0.360 to 0.690), with one seed collapsing on its phase path\\
Learned warm start for phase-fixed WGS (refuted) & \footnotesize measured negative: on held-out geometries the network start needs a median 8 iterations to 1 percent against 14 from a random start (1.75 times, below the 2 times claimed), and a control with the same early phase fix and no network gains as much; the model is not released and the claim learned warm start is reported as refuted\\
Aberration-robust hologram & \footnotesize over a 200-sample ensemble of residual aberrations, the 95th percentile of the relative standard deviation at least 1.25 times below phase-fixed weighted GS, efficiency at most 10 percent lower (measured: a median of 1.435 times, 1.42 to 1.57 over three seeds, at 0.982 of the efficiency); robustness to pointing errors is measured and reported, not claimed (1.04 times at the acceptance budget; 1.48 times at 0.10 bins and 0.912 of the efficiency with a displacement prior alone and five times the budget)\\
Efficiency-uniformity Pareto front & \footnotesize hypervolume at least that of the GS, generalized adaptive additive GS, direct search and weighted GS points on the 10x10 task (measured: 0.8190 against 0.8018); in the benchmark the stored front (efficiency-uniformity Pareto front's converged run of the first seed) reaches that hypervolume on none of the 6 geometries with a front against the medians at 100 iterations on the ideal device, below them by at most 0.007, and on all 6 against those at 30 (the count at 100 iterations moves between runs of the benchmark, 3, 4 and 0, because a cell's medians depend on the seeds its wall-clock bound admits), 2026-10-31 run\\
\bottomrule\end{longtable}

\begin{longtable}{@{}>{\raggedright\arraybackslash}p{0.2\textwidth} p{0.58\textwidth} p{0.14\textwidth}@{}}
\caption{Measured acceptance of the methods on DinHot's own grids (unit tests; seeds as stated). Intervals are [min, max].}\label{tab:acceptance}\\
\toprule Method & Measured & Seeds\\\midrule\endfirsthead
\toprule Method & Measured & Seeds\\\midrule\endhead
\footnotesize Superposition of prisms and lenses & \footnotesize efficiency 0.281, uniformity 0.0003 (deterministic) & \footnotesize 1\\
\footnotesize Random superposition & \footnotesize efficiency 0.750 [0.738, 0.752], uniformity 0.252 [0.106, 0.272] & \footnotesize 6\\
\footnotesize Random mask & \footnotesize 256 squared: efficiency 0.010, uniformity 0.057 [0.031, 0.072]; 768 squared: efficiency 0.008, uniformity 0.566 [0.543, 0.590] & \footnotesize 5 and 2\\
\footnotesize Gerchberg-Saxton & \footnotesize flat beam: efficiency 0.939 [0.933, 0.939], uniformity 0.604 [0.557, 0.626]; Gaussian beam: efficiency 0.815 against its ceiling 0.869 & \footnotesize 3\\
\footnotesize Generalized adaptive additive GS & \footnotesize uniformity 0.688 [0.638, 0.709] against 0.615 [0.550, 0.637] for GS at equal iterations; efficiency 0.811 against 0.816 & \footnotesize 5\\
\footnotesize Weighted GS & \footnotesize flat beam: efficiency and uniformity both at 0.9 after 12, 15 and 12 iterations (median uniformity 0.857 at 10); Gaussian beam: 12, 15 and 12; 3 x 3 x 3 lattice (0.9 and 0.98): 35, 20 and 27; budgets 15, 20 and 100 & \footnotesize 3\\
\footnotesize Phase-fixed weighted GS & \footnotesize 50 x 30 array to 0.5 percent in 64 and 64 iterations against 136 and 106 for weighted GS (phase fixed after 60; efficiency 0.83 and 0.83); with the source value of 12 it stalls at 3.58 percent & \footnotesize 2\\
\footnotesize Mixed-region amplitude freedom & \footnotesize signal-region error 12.4 (disc), 7.8 (ring) and 10.4 (bar) times below GS at the worse of two starts; efficiency 0.181 to 0.184 against 0.963 to 0.992 & \footnotesize 2 starts\\
\footnotesize Direct search & \footnotesize uniformity 0.9994 and 0.9992 (efficiency 0.66 and 0.65); 119 quasi-periodic traps at 0.023 percent rms & \footnotesize 2 and 1\\
\footnotesize Conjugate-gradient minimization & \footnotesize analytic gradient within 2.5e-06 of finite differences; square lattice error 0.55 and 0.51 percent & \footnotesize 2\\
\footnotesize Wirtinger gradient descent with Adam & \footnotesize relative standard deviation 2.53, 2.46 and 2.26 percent against 17.78, 17.98 and 18.82 for GS at 30 iterations & \footnotesize 3\\
\footnotesize Closed-loop weighted GS & \footnotesize emulated camera loop: open-loop 20.9 and 20.4 percent to 0.84 and 1.07 percent in 20 iterations; 24.3 and 23.6 percent to 1.05 and 0.74 percent in three phase-fixed steps & \footnotesize 2\\
\footnotesize Device-aware CG with pixel crosstalk & \footnotesize error 0.241 (efficiency 0.535) against 1.382 (0.507) for conjugate-gradient minimization without the device model, large disc on the crosstalk device & \footnotesize 1\\
\footnotesize Zernike correction and sensorless modal recovery & \footnotesize worst per-mode relative error 0.002 [0.001, 0.002] over Noll modes 4 to 11 & \footnotesize 3\\
\footnotesize Laguerre-Gaussian vortex traps & \footnotesize ring phase winds by 2 pi l (winding error below the 1e-9 rad of the arithmetic) with a maximum error of 0.033 [0.028, 0.060] rad (isolated vortices of charge 1, 2, 3 and -2) & \footnotesize 2 per charge\\
\footnotesize Hungarian assignment and linear phase interpolation & \footnotesize 0 path crossings in 60 loadings; moving traps keep 0.767 [0.747, 0.792] of their mean intensity (worst moving trap of each of 5 loadings) & \footnotesize 60 loadings\\
\footnotesize Multi-plane CNN hologram & \footnotesize three-plane network refined by volumetric 3D GS: median accuracy ratio 1.140 (quartiles 1.103 and 1.163) on 162 held-out volumes; taken from the committed metrics.json, not remeasured here & \footnotesize 162 volumes\\
\footnotesize Optimal-transport initialized MRAF & \footnotesize signal-region error 0.0204 (disc), 0.0179 (ring) and 0.0221 (line) against 0.250, 0.188 and 0.223 for MRAF from a random start (the better of two); baked size: 0.0259 (disc) and 0.0171 (ring) against 0.266 and 0.228 & \footnotesize 2 and 1\\
\footnotesize Volumetric 3D GS & \footnotesize plane means over overall mean 0.938, 0.987 and 1.075 and 0.926, 1.034 and 1.039 on three planes & \footnotesize 2\\
\footnotesize Device-aware weighted GS & \footnotesize crosstalk device, 10 x 10 task: relative standard deviation 0.20 [0.16, 0.30] percent against 5.21 [4.75, 5.46] for phase-fixed weighted GS at 100 iterations (a factor 26 [16, 33]) at 0.995 [0.991, 0.997] of its efficiency; 4-bit device with fill 0.85: 0.72 [0.60, 1.16] against 5.73 [5.13, 5.88] (a factor 7.9 [5.0, 8.6]) at 0.990 [0.983, 0.996] of its efficiency; taken from the companion manuscript's beyond.json, 5 and 3 seeds & \footnotesize 5 and 3\\
\footnotesize Flicker-bounded phase-continuous sequence & \footnotesize 4 x 4 lattice, two columns moved by four bins in 13 frames: lowest transient intensity of a moving trap 0.920 [0.915, 0.932] of its mean against 0.654 [0.551, 0.709] for linear phase interpolation; largest per-frame spread 0.34 [0.14, 0.39] against 22.8 [19.2, 27.0] percent; taken from the companion manuscript's beyond.json, 5 seeds & \footnotesize 5\\
\footnotesize Aberration- and pointing-robust hologram & \footnotesize 95th percentile of the relative standard deviation a median 1.435 times below phase-fixed weighted GS (1.42 to 1.57) over a 200-sample aberration ensemble, at a median 0.982 of its efficiency; pointing errors 1.04 times (seed 1) & \footnotesize 3 (pointing 1)\\
\footnotesize Efficiency-uniformity Pareto front & \footnotesize hypervolume 0.8190 against 0.8018 for the union of the GS, generalized adaptive additive GS, weighted GS and direct search points; 7 front points dominate GS and 8 dominate generalized adaptive additive GS & \footnotesize 1\\
\bottomrule\end{longtable}


\subsection{Seeds, repetitions, hardware and run times}\label{sec:protocol-hw}

Every method of a case receives the same seed, so that all of them start from the same random numbers. The acceptance tests state their seeds in Table~\ref{tab:acceptance}: between one and six, the number a test needs to bound a measurement within the time of the test suite. The expected ranges of the cases are measured on three seeds. The benchmark repeats the methods on its geometries over up to \benchSeeds{} seeds; a group runs seeds until a bound on its wall clock, and at least two, so a cell carries between \benchSeedsMin{} and \benchSeedsMax{} seeds, which the file records per row. The uncertainty reported is the observed range over the seeds for the acceptance tests and the bake, and the median with the inter-quartile range for the benchmark; no parametric interval is claimed, because the seeds are few and the distributions are not Gaussian.

The measuring machine is the development workstation recorded in the product's file of measured resources: a laptop processor with twelve cores, forty-eight gigabytes of memory and a laptop graphics processor, Python 3.13 and Node 24, with the browser measurements taken in Chromium driven by Playwright. The canonical bake of the eleven cases that run in parallel took \bakeIdleS{} s on an idle machine and \bakeLoadedS{} s while the machine ran other jobs, at a peak of \bakeCaseMB{} MB for one case process; the benchmark and the training are bounded by the design at hours and run locally, and the times in the benchmark are wall-clock seconds of the machine that ran it, under shared load, so they differ between runs and are used only to compare methods inside one run. The browser timings of Section~\ref{sec:live} are quoted from the same record.

\subsection{The benchmark protocol}\label{sec:protocol-bench}

The benchmark compares the methods on one protocol. Its geometries are (\benchGeometries{} in all) a $10\times10$ square, a ring of 36 traps, 100 random traps with a minimum spacing, a quasi-periodic set of 119 traps, a $5\times5\times3$ lattice, a large set (the $32\times32$ lattice and the disordered 819-trap set, pooled by seed parity) and four shaped fields (a ring, a disc, a line and a Laguerre--Gaussian ring, pooled by seed). The grids are \benchGrids{} pixels, with the benchmark optics keeping the Fourier unit and the focal bin fixed across grids; the iteration budgets are \benchIters{} and the time budgets are \benchTimes{} s, with a converged row run to a cap of iterations. For every method, geometry, device and parameter set and for every seed, six runs are made; every cell reports the median and the inter-quartile range of efficiency, uniformity, relative standard deviation, time and iterations, on the ideal device and on the C08 device (every hologram scored against its device). The device rows run with the default parameters only and not on the multi-plane geometries.

A method that declares a tuning grid of four candidates, or has a fallback entry whose names its run function accepts, is tuned: each candidate runs on three other seeds at 30 iterations and the one with the lowest median relative standard deviation among those with at least half the best median efficiency gives the tuned rows. The number of tuning runs is the same for every tuned method, and no row mixes one method's tuned parameters with another's defaults, as the parameter field of each row states. The committed results file holds \benchRows{} rows of \benchNMeasured{} of the \benchNProtocol{} engine methods; \benchSkipped{} method--geometry cells are recorded as skipped with the reason the file states (a method that does not accept the target, such as a two-dimensional method on a three-dimensional lattice, or one that refuses the geometry), and the \benchNNoRows{} methods without a row (\benchNoRows) refuse every geometry of the protocol by their specification (a vortex charge, a rearrangement block, the geometry the learned model was trained for); their tests and their cases measure them.

\subsection{Evaluation of the learned models}\label{sec:protocol-learned}

The two learned methods are evaluated only on held-out geometries that the training never saw. The held-out set is written before training and its hash is recorded in the model card. The metrics document compares each model with its baseline: the volumetric GS at the same wall-clock budget for the network (the budget is the median time of the baseline at 30 iterations on those volumes), and the random start for the warm start, with the same early freeze as a control without the network. A model that does not meet its criterion is not released: the baseline answers, readiness says training, and the measured values are reported as they are. The rules that keep the evaluation sound are stated before the results: held-out means disjoint from calibration and training; the acceptance values of each method state seeds, grid and date; accurate means that the result agrees with the stated engine and device model, not that it matches a real instrument; and a published number is labelled with its source and conditions while a DinHot measurement is labelled as one.
